% Part: second-order-logic
% Chapter: syntax-and-semantics

\documentclass[../../../include/open-logic-chapter]{subfiles}

\begin{document}

\olchapter{sol}{syn}{সংকেতবিন্যাস ও অর্থতত্ত্ব}

\begin{editorial}
এ পর্যন্ত SOL-এর মৌলিক সংকেতবিন্যাস ও অর্থতত্ত্ব আলোচনা করা হয়েছে।
অধ্যায় হিসেবে এটি খুবই ছোট। SOL-এর !!{derivation} ব্যবস্থা নিয়ে কথা
বলতে হলে দ্বিতীয়-ক্রমীয় চলরাশির প্রতিস্থাপনও আলোচনা করতে হবে, আর
সেখানে কয়েকটি সূক্ষ্ম সমস্যা আছে।
\end{editorial}

\olimport{introduction}

\olimport{terms-formulas}

\olimport{satisfaction}

\olimport{semantic-notions}

\olimport{expressive-power}

\olimport{inf-count}

\OLEndChapterHook

\end{document}
